% !TEX root = ../../../main.tex
% 所属：第四部分：Fourier 分析 / 第4章：$\R$上的Fourier变换
% 节：4.4 Poisson 公式
% 正文按原文件顺序保存；由 ../chapter4.tex 调用。
\section{Poisson 公式}

\begin{definition}
	 给定实直线上的函数 $f\in\mathcal S(\R)$，定义
\[
F_1(x)=\sum_{n=-\infty}^{\infty}f(x+n)
\]
\end{definition}
得到一个定义在圆周上的函数。由于 $f$ 是迅速下降的，这个级数在 $\R$ 的任意紧集上一致且绝对收敛，因此 $F_1$ 连续。将求和指标 $n$ 替换为 $n+1$，可得
\[
F_1(x+1)=F_1(x),
\]
所以 $F_1$ 是以 $1$ 为周期的函数，称 $F_1$ 为函数 $f$ 的\textbf{周期化}。

\begin{definition}
	由 Fourier 反演公式，
\[
f(x)=\int_{-\infty}^{\infty}\widehat f(\xi)e^{2\pi i\xi x}\dd\xi.
\]
考虑该式的离散化形式，即用求和代替积分，得到
\[
F_2(x)=\sum_{n=-\infty}^{\infty}\widehat f(n)e^{2\pi i n x}.
\]
\end{definition}
由于 $\widehat f\in\mathcal S(\R)$，这个级数一致且绝对收敛，因此 $F_2$ 连续。又因为每个指数函数 $e^{2\pi i n x}$ 都以 $1$ 为周期，故 $F_2$ 也是以 $1$ 为周期的函数。

\begin{theorem}
	\textbf{（Poisson 求和公式）}
	若 $f\in\mathcal S(\R)$，则
	\[
	\sum_{n=-\infty}^{\infty}f(x+n)
	=
	\sum_{n=-\infty}^{\infty}\widehat f(n)e^{2\pi i n x}.
	\]
	特别地，取 $x=0$，有
	\[
	\sum_{n=-\infty}^{\infty}f(n)
	=
	\sum_{n=-\infty}^{\infty}\widehat f(n).
	\]
	\label{thm:poisson-summation}
\end{theorem}
\newpage
\begin{proof}
	由\ref{Fourier级数的唯一性},只需证明等式两边作为连续的 $1$-周期函数，具有相同的 Fourier 系数。右端的第 $m$ 项 Fourier 系数显然为 $\widehat f(m)$。对于左端，由局部一致绝对收敛可交换求和与积分，故
	\[
	\begin{aligned}
		&\int_0^1\left(\sum_{n=-\infty}^{\infty}f(x+n)\right)e^{-2\pi i m x}\dd x=\sum_{n=-\infty}^{\infty}\int_0^1f(x+n)e^{-2\pi i m x}\dd x\\
		&\qquad=\sum_{n=-\infty}^{\infty}\int_n^{n+1}f(y)e^{-2\pi i m y}\dd y=\int_{-\infty}^{\infty}f(y)e^{-2\pi i m y}\dd y
		=\widehat f(m).
	\end{aligned}
	\]
	由连续周期函数的 Fourier 级数唯一性，等式成立。
\end{proof}

\begin{definition}
	 当 $s>0$ 时，定义 Theta 函数
\[
\vartheta(s)=\sum_{n=-\infty}^{\infty}e^{-\pi n^2s}.
\]
\end{definition}
$\vartheta$函数满足这个重要的函数方程。
\begin{theorem}	\label{thm:theta-functional-equation}
	当 $s>0$ 时，有
	\[
	s^{-1/2}\vartheta(1/s)=\vartheta(s).
	\]
\end{theorem}

\begin{proof}
	取
	\[
	f(x)=e^{-\pi s x^2}.
	\]
	由 Gaussian 函数的 Fourier 变换公式，
	\[
	\widehat f(\xi)=s^{-1/2}e^{-\pi\xi^2/s}.
	\]
	在 Poisson 求和公式中取 $x=0$，即得
	\[
	\vartheta(s)
	=\sum_{n=-\infty}^{\infty}e^{-\pi n^2s}
	=s^{-1/2}\sum_{n=-\infty}^{\infty}e^{-\pi n^2/s}
	=s^{-1/2}\vartheta(1/s).
	\]
\end{proof}

Theta 函数 $\vartheta(s)$ 还可以拓展至复数域 $\operatorname{Re}(s)>0$，并且上面的函数方程仍然成立。在数论中，Theta 函数与定义在 $\operatorname{Re}(s)>1$ 上的 Riemann $\zeta$ 函数
\[
\zeta(s)=\sum_{n=1}^{\infty}\frac1{n^s}
\]
有着本质联系。以后会看到，这个函数包含了素数的一些基本信息（见第 7 章）。

\begin{proposition}	\label{prop:theta-zeta-mellin}
	当 $s>1$ 时，
	\[
	\pi^{-s/2}\Gamma\left(\frac{s}{2}\right)\zeta(s)
	=
	\frac12\int_0^{\infty}t^{s/2-1}\bigl[\vartheta(t)-1\bigr]\dd t.
	\]
\end{proposition}

\begin{proof}
	由
	\[
	\vartheta(t)
	=\sum_{n=-\infty}^{\infty}e^{-\pi n^2t}
	=1+2\sum_{n=1}^{\infty}e^{-\pi n^2t},
	\]
	可得
	\[
	\frac12\bigl[\vartheta(t)-1\bigr]
	=\sum_{n=1}^{\infty}e^{-\pi n^2t}.
	\]
	令
	\[
	F_N(t)=\sum_{n=1}^{N}t^{s/2-1}e^{-\pi n^2t}.
	\]
	则 $F_N(t)$ 单调增加到 $t^{s/2-1}\sum\limits_{n=1}^{\infty}e^{-\pi n^2t}$。由单调收敛定理，可交换求和与积分：
	\[
	\frac12\int_0^{\infty}t^{s/2-1}\bigl[\vartheta(t)-1\bigr]\dd t
	=
	\sum_{n=1}^{\infty}\int_0^{\infty}t^{s/2-1}e^{-\pi n^2t}\dd t.
	\]
	令 $u=\pi n^2t$，则
	\[
	\begin{aligned}
		\int_0^{\infty}t^{s/2-1}e^{-\pi n^2t}\dd t
		&=(\pi n^2)^{-s/2}
		\int_0^{\infty}u^{s/2-1}e^{-u}\dd u\\
		&=\pi^{-s/2}n^{-s}\Gamma\left(\frac{s}{2}\right).
	\end{aligned}
	\]
	因此
	\[
	\frac12\int_0^{\infty}t^{s/2-1}\bigl[\vartheta(t)-1\bigr]\dd t
	=
	\pi^{-s/2}\Gamma\left(\frac{s}{2}\right)
	\sum_{n=1}^{\infty}\frac1{n^s},
	\]
	即得所证。条件 $s>1$ 保证 $\zeta(s)$ 收敛。
\end{proof}

\begin{definition}
定义一般化的 Theta 函数
	\[
	\Theta(z\mid\tau)
	=
	\sum_{n=-\infty}^{\infty}
	e^{i\pi n^2\tau}e^{2\pi i n z},
	\]
	其中 $\operatorname{Im}(\tau)>0$ 且 $z\in\C$。此时级数绝对收敛。取 $z=0$ 以及 $\tau=is$，便得到
	\[
	\Theta(0\mid is)=\vartheta(s).
	\]
\end{definition}
\newpage
下面讨论圆盘上和上半平面上的 Poisson 核，即
\[
P_r(\theta)=\frac{1-r^2}{1-2r\cos\theta+r^2},
\qquad
P_y(x)=\frac1\pi\frac{y}{y^2+x^2},
\]
以及它们之间的联系。

\begin{theorem}
	当 $r=e^{-2\pi y}$ 时，有
	\[
	P_r(2\pi x)=\sum_{n\in\Z}P_y(x+n).
	\]
	\label{thm:poisson-kernel-periodization}
\end{theorem}

\begin{proof}
	由标准 Fourier 变换公式
	\[
	\widehat{P_y}(\xi)=e^{-2\pi y\abs{\xi}}.
	\]
	由于 $P_y(x+n)=O(n^{-2})$，级数 $\sum_{n\in\Z}P_y(x+n)$ 在 $[0,1]$ 上一致且绝对收敛，因而其和连续且以 $1$ 为周期。它的第 $m$ 个 Fourier 系数为
	\[
	\int_{-\infty}^{\infty}P_y(x)e^{-2\pi i m x}\dd x
	=e^{-2\pi y\abs{m}}
	=r^{\abs{m}}.
	\]
	另一方面，圆盘 Poisson 核的 Fourier 展开为
	\[
	P_r(2\pi x)=\sum_{m\in\Z}r^{\abs{m}}e^{2\pi i m x}.
	\]
	两边具有相同的 Fourier 系数，故由唯一性定理得到结论。
\end{proof}
